\section{Детерминизм и «квантовая случайность''}

\begin{proposition}[Детерминизм $\Mlayer$]
\label{prop:determinism}
Поскольку новое состояние есть функция предыдущего,
\[
  \symCapitalPsi^{t+1}=g(\symCapitalPsi^t),
\]
на каждом такте значение одно, и траектория $\symCapitalPsi^t$ восстанавливается из $\symCapitalPsi^0$.
\end{proposition}

\begin{proof}
\pfrom{определения $\symCapitalPsi^{t+1}=g(\symCapitalPsi^t)$}
при фиксированном $\symCapitalPsi^t$ правая часть однозначна.
\phence индукция по~$t$ восстанавливает $\symCapitalPsi^t$ из $\symCapitalPsi^0$ и правила~$g$.
\end{proof}

На $\Tlayer$ частота исхода есть квадрат модуля амплитуды \eqref{eq:born};
она возникает из динамического хаоса фаз и грубости прибора (\cref{fig:axiom-layer-m,fig:axiom-layer-t}).

\begin{figure}[htbp]
  \centering
  \begin{tikzpicture}[carrier panel, x=0.32cm, y=0.32cm]
  \foreach \r in {0,...,15} {
    \foreach \c in {0,...,15} {
      \pgfmathsetmacro{\v}{cos(12*\r + 7*\c)}
      \pgfmathsetmacro{\g}{int(50 + 40*\v)}
      \fill[gray!\g] (\c,15-\r) rectangle ++(1,1);
    }
  }
\end{tikzpicture}

  \caption{$\mathcal{M}$: детерминированное микрополе $z(x,t)$.}
  \label{fig:axiom-layer-m}
\end{figure}

\begin{figure}[htbp]
  \centering
  \begin{tikzpicture}[carrier panel, x=0.32cm, y=0.32cm]
  \foreach \r in {0,...,15} {
    \foreach \c in {0,...,15} {
      \pgfmathsetmacro{\v}{0.6*cos(12*\r + 7*\c) + 0.25*cos(3*\r) + 0.2*cos(3*\c)}
      \pgfmathsetmacro{\g}{int(50 + 35*\v)}
      \fill[gray!\g] (\c,15-\r) rectangle ++(1,1);
    }
  }
\end{tikzpicture}

  \caption{$\mathcal{T}$: грубый макро-readout $\Phi=\mathcal{B}z$.}
  \label{fig:axiom-layer-t}
\end{figure}

